% Original timeline; winding path inspired by user reference, with no decorative icons.
\begin{tikzpicture}[x=1mm,y=1mm,font=\figurefont\footnotesize,
 milestone/.style={circle,minimum size=7mm,inner sep=0pt,draw=BookTeal,fill=BookTeal!20,line width=.8pt},
 title/.style={font=\figurefont\small,text=BookInk},
 detail/.style={font=\figurefont\footnotesize,text=BookMuted}]
\path[use as bounding box] (0,-8) rectangle (155,182);
\draw[BookMuted!60,line width=3.5mm,line cap=round]
 (78,178) .. controls (48,174) and (48,164) .. (78,158)
 .. controls (108,152) and (108,142) .. (78,136)
 .. controls (48,130) and (48,120) .. (78,114)
 .. controls (108,108) and (108,98) .. (78,92)
 .. controls (48,86) and (48,76) .. (78,70)
 .. controls (108,64) and (108,54) .. (78,48)
 .. controls (48,42) and (48,32) .. (78,26)
 .. controls (108,20) and (108,10) .. (78,4);
\draw[white,dashed,line width=.5mm,dash pattern=on 2mm off 2mm]
 (78,178) .. controls (48,174) and (48,164) .. (78,158)
 .. controls (108,152) and (108,142) .. (78,136)
 .. controls (48,130) and (48,120) .. (78,114)
 .. controls (108,108) and (108,98) .. (78,92)
 .. controls (48,86) and (48,76) .. (78,70)
 .. controls (108,64) and (108,54) .. (78,48)
 .. controls (48,42) and (48,32) .. (78,26)
 .. controls (108,20) and (108,10) .. (78,4);
\node[milestone] at (56,170) {};
\node[anchor=west,text=BookTeal,font=\figurefont\small] at (2,175) {1943};
\node[title,anchor=west,text width=49mm,align=left] at (2,169) {阈值神经元};
\node[detail,anchor=north west,text width=49mm,align=left] at (2,165) {把输入组合成离散响应};
\node[milestone] at (100,148) {};
\node[anchor=west,text=BookTeal,font=\figurefont\small] at (109,153) {1958};
\node[title,anchor=west,text width=46mm,align=left] at (109,147) {感知机};
\node[detail,anchor=north west,text width=46mm,align=left] at (109,143) {由分类错误更新线性分界};
\node[milestone] at (56,126) {};
\node[anchor=west,text=BookTeal,font=\figurefont\small] at (2,131) {1986};
\node[title,anchor=west,text width=49mm,align=left] at (2,125) {反向传播};
\node[detail,anchor=north west,text width=49mm,align=left] at (2,121) {沿复合映射分配参数梯度};
\node[milestone] at (100,104) {};
\node[anchor=west,text=BookTeal,font=\figurefont\small] at (109,109) {1997—1998};
\node[title,anchor=west,text width=46mm,align=left] at (109,103) {LSTM与LeNet};
\node[detail,anchor=north west,text width=46mm,align=left] at (109,99) {以门控记忆和局部共享形成结构};
\node[milestone] at (56,82) {};
\node[anchor=west,text=BookTeal,font=\figurefont\small] at (2,87) {2012};
\node[title,anchor=west,text width=49mm,align=left] at (2,81) {AlexNet};
\node[detail,anchor=north west,text width=49mm,align=left] at (2,77) {数据与计算支持深层视觉学习};
\node[milestone] at (100,60) {};
\node[anchor=west,text=BookTeal,font=\figurefont\small] at (109,65) {2017};
\node[title,anchor=west,text width=46mm,align=left] at (109,59) {Transformer};
\node[detail,anchor=north west,text width=46mm,align=left] at (109,55) {以注意力组织跨位置读取};
\node[milestone] at (56,38) {};
\node[anchor=west,text=BookTeal,font=\figurefont\small] at (2,43) {2020};
\node[title,anchor=west,text width=49mm,align=left] at (2,37) {规模与自监督};
\node[detail,anchor=north west,text width=49mm,align=left] at (2,33) {以大规模预训练学习通用表示};
\node[milestone] at (100,16) {};
\node[anchor=west,text=BookTeal,font=\figurefont\small] at (109,21) {2024};
\node[title,anchor=west,text width=46mm,align=left] at (109,15) {状态空间与条件计算};
\node[detail,anchor=north west,text width=46mm,align=left] at (109,11) {研究长序列状态、专家路由及效率};
\end{tikzpicture}
